\chapter*{Kata Pengantar}
\addcontentsline{toc}{chapter}{Kata Pengantar}

Puji syukur kami panjatkan ke hadirat Tuhan Yang Maha Esa atas segala rahmat dan karunia-Nya sehingga penulisan "Modul Praktikum Prediksi Awal Kesehatan Daun Berbasis Citra Digital" ini dapat diselesaikan dengan baik.

Modul ini dirancang secara khusus untuk memberikan pengalaman praktis secara komprehensif dari awal hingga akhir (end-to-end) dalam membangun sebuah sistem Machine Learning berbasis Computer Vision. Melalui studi kasus klasifikasi kondisi daun menjadi 5 (lima) kelas, pembaca akan dipandu mulai dari tahapan pemuatan data, eksplorasi data, perancangan model jaringan saraf tiruan (ResNet9), pelatihan model, hingga tahapan \textit{deployment} menjadi sebuah aplikasi web sederhana menggunakan Streamlit.

Penggunaan Google Colab sebagai environment utama dalam praktikum ini diharapkan dapat memudahkan pembaca dalam mempraktekkan materi tanpa perlu memikirkan spesifikasi perangkat keras yang tinggi, karena seluruh komputasi dilakukan melalui teknologi \textit{cloud}.

Kami menyadari bahwa modul ini tidak luput dari kekurangan. Oleh karena itu, segala bentuk kritik dan saran yang membangun sangat kami harapkan guna penyempurnaan di masa yang akan datang. Perlu ditegaskan pula bahwa model yang dibangun dalam modul ini berfungsi sebagai purwarupa (prototype) untuk \textit{screening} awal, dan bukan sebagai diagnosis agrikultural yang final.

Akhir kata, semoga modul ini dapat memberikan manfaat yang sebesar-besarnya bagi para pembaca, khususnya bagi pihak yang sedang mendalami implementasi Machine Learning dan Computer Vision di dunia nyata.

\vspace{2cm}
\begin{flushright}
Penyusun
\end{flushright}
